\chapter{Annexes}
\label{ch:annexes}

Les annexes (40 pages maximum) rassemblent les éléments de la fonctionnalité la plus
représentative : maquettes, captures et code, code des composants métier et d'accès aux données,
code des autres composants, et jeux de tests complets.

\section{Gestion de projet : suivi et comptes rendus}
\label{ann:gestion}

\subsection{Extrait du fichier de suivi}

Le fichier \file{SUIVI.md}, consolidé depuis l'historique Git, recense les tâches du projet par
phase (tâche, phase, date). Le projet s'est déroulé en deux temps : une construction initiale
(2024), puis une reprise dédiée à la sécurisation et à l'industrialisation (2026).

\paragraph{Phase 1 — Construction initiale (sept.–nov. 2024).}
\begin{itemize}\setlength{\itemsep}{0pt}
  \item page d'accueil et fil d'actualité (front) ;
  \item « j'aime » (likes) et appels AJAX ;
  \item routeur et gestion des utilisateurs \emph{(amorce contribuée par un binôme, depuis
        entièrement reprise)} ;
  \item module de réponses imbriquées, navigation et routes ;
  \item minification des assets, comptage récursif des réponses.
\end{itemize}

\paragraph{Phase 2 — Reprise : sécurisation \& industrialisation (mai–juin 2026).}
\begin{itemize}\setlength{\itemsep}{0pt}
  \item conteneurisation (Docker) et migration de la base vers PostgreSQL ;
  \item gestion des secrets (\file{.env}) et logger PSR-3 ;
  \item \texttt{AdminMiddleware} et protection des routes d'administration ;
  \item échappement anti-XSS et contrôle d'accès anti-IDOR sur la suppression ;
  \item refonte MVC (correction du N+1, découplage du validateur, unification JSON) ;
  \item protection CSRF, en-têtes de sécurité, injection de dépendance de session ;
  \item jeu d'essai (\file{seed.sql}), PHP\_CodeSniffer (PSR-12), refonte de la feuille de style.
\end{itemize}

\subsection{Comptes rendus de session}

Reconstitués depuis l'historique Git et le journal de décisions (\file{Probleme-Solution.md}).
Chaque session est rattachée à un ensemble de commits datés.

\paragraph{6 mai 2026 — Industrialisation de l'environnement.}
\textbf{Objectif} : rendre l'environnement reproductible et conforme à la production.
\textbf{Réalisé} : conteneurisation (FrankenPHP, Caddy, Traefik) ; migration de MariaDB vers
PostgreSQL ; sortie des identifiants de connexion vers \file{.env} (phpdotenv) ; logger PSR-3 ;
ajout d'\texttt{AdminMiddleware} et durcissement de la gestion d'erreurs.
\textbf{Reste à faire} : valider le type MIME réel des téléversements ; régénérer l'identifiant
de session après connexion (anti-fixation).

\paragraph{12 mai 2026 — Audit MVC et failles applicatives.}
\textbf{Objectif} : remettre la couche présentation à sa place et corriger les failles
identifiées. \textbf{Réalisé} : échappement systématique du contenu utilisateur (anti-XSS) ;
contrôle d'accès « propriétaire ou administrateur » sur la suppression (anti-IDOR) ; préchargement
des auteurs en une requête (correction du N+1, 21 requêtes ramenées à 2 sur le fil) ; extraction
du validateur dans \file{modules/validators/} ; suppression de \texttt{DynamicMessageController}
et unification du format de réponse JSON ; correction de quatre bugs préexistants.
\textbf{Reste à faire} : doubler systématiquement la validation côté serveur (format e-mail,
robustesse du mot de passe) ; mettre en place des tests automatisés.

\paragraph{13 mai 2026 — Jeu d'essai et charte graphique.}
\textbf{Objectif} : disposer d'un jeu de données reproductible et d'une interface cohérente.
\textbf{Réalisé} : extraction des données de démonstration dans \file{seed.sql} (utilisateurs,
messages, « j'aime », réponses) ; refonte de la feuille de style.
\textbf{Reste à faire} : documenter et tester la procédure de sauvegarde / restauration de la base
(\texttt{pg\_dump} / \texttt{pg\_restore}).

\paragraph{15 juin 2026 — Durcissement de la sécurité et outillage qualité.}
\textbf{Objectif} : compléter la défense en profondeur et industrialiser le contrôle de qualité.
\textbf{Réalisé} : protection CSRF par jeton de synchronisation armé sur tout POST ; en-têtes de
sécurité (CSP, \texttt{nosniff}, HSTS) ; injection de dépendance du gestionnaire de session ; mise
en place de PHP\_CodeSniffer (PSR-12) et reformatage de la base de code.
\textbf{Reste à faire} : automatiser l'analyse statique et les tests via une chaîne d'intégration
continue ; constituer une couverture de tests (PHPUnit).

\section{Script de création de la base de données}
\label{ann:sql}

Script complet \file{solage.pg.sql}, chargé à l'initialisation du conteneur PostgreSQL.

\begin{lstlisting}[style=sql]
BEGIN;

CREATE TABLE roles (
    id   SERIAL PRIMARY KEY,
    name VARCHAR(255) NOT NULL
);
INSERT INTO roles (id, name) VALUES (1,'Admin'), (2,'Utilisateur'), (3,'Moderateur');

CREATE TABLE users (
    id        SERIAL PRIMARY KEY,
    name      VARCHAR(255) NOT NULL,
    firstname VARCHAR(255) NULL,
    email     VARCHAR(255) NOT NULL UNIQUE,
    password  VARCHAR(255) NOT NULL,
    role      INT NULL REFERENCES roles(id) ON DELETE SET NULL,
    image     VARCHAR(255) NULL
);
CREATE INDEX users_role_idx ON users(role);

CREATE TABLE posts (
    id              SERIAL PRIMARY KEY,
    user_id         INT NULL REFERENCES users(id) ON DELETE CASCADE,
    date            TIMESTAMP NOT NULL,
    likes           INT DEFAULT 0,
    content         TEXT NULL,
    reply_to        INT NULL,
    reply_to_parent INT NULL,
    image           TEXT NULL
);
CREATE INDEX posts_user_idx ON posts(user_id);

CREATE TABLE responses (
    id       SERIAL PRIMARY KEY,
    post     INT NULL REFERENCES posts(id) ON DELETE CASCADE,
    user_id  INT NULL REFERENCES users(id) ON DELETE CASCADE,
    date     TIMESTAMP NOT NULL,
    likes    INT DEFAULT 0,
    content  TEXT NULL,
    reply_to INT NULL REFERENCES responses(id) ON DELETE SET NULL
);
CREATE INDEX responses_post_idx     ON responses(post);
CREATE INDEX responses_user_idx     ON responses(user_id);
CREATE INDEX responses_reply_to_idx ON responses(reply_to);

CREATE TABLE likes (
    id         SERIAL PRIMARY KEY,
    user_id    INT NULL REFERENCES users(id) ON DELETE CASCADE,
    post       INT NULL REFERENCES posts(id) ON DELETE CASCADE,
    response   INT NULL REFERENCES responses(id) ON DELETE CASCADE,
    created_at TIMESTAMP NOT NULL
);
CREATE INDEX likes_user_idx     ON likes(user_id);
CREATE INDEX likes_post_idx     ON likes(post);
CREATE INDEX likes_response_idx ON likes(response);

CREATE TABLE users_favorites_posts (
    user_id INT NOT NULL REFERENCES users(id) ON DELETE CASCADE,
    post    INT NOT NULL REFERENCES posts(id) ON DELETE CASCADE,
    PRIMARY KEY (user_id, post)
);
CREATE INDEX users_favorites_posts_post_idx ON users_favorites_posts(post);

COMMIT;
\end{lstlisting}

\section{Code d'un autre composant : le routeur}
\label{ann:router}

Routeur complet \file{src/Router.php} : correspondance d'URL, exécution des middlewares, armement
du CSRF sur les requêtes POST, et instanciation du contrôleur.

\begin{lstlisting}[style=php]
class Router {
    protected $routes = [];

    public function addRoute($method, $path, $target, $middleware = null) {
        $this->routes[] = ['method' => $method, 'path' => $path,
                           'target' => $target, 'middleware' => $middleware];
    }

    public function match() {
        $requestUri    = parse_url($_SERVER['REQUEST_URI'], PHP_URL_PATH);
        $requestMethod = $_SERVER['REQUEST_METHOD'];

        foreach ($this->routes as $route) {
            $pathRegex = preg_replace('/\{([a-zA-Z0-9_]+)\}/', '([a-zA-Z0-9_]+)', $route['path']);
            $pathRegex = '#^' . $pathRegex . '$#';

            if ($route['method'] === $requestMethod && preg_match($pathRegex, $requestUri, $matches)) {
                if ($route['middleware']) {
                    (new $route['middleware']())->handle();
                }
                if ($route['method'] === 'POST') {
                    (new CsrfMiddleware())->handle();
                }
                $routeArray = explode('#', $route['target']);
                if (count($routeArray) < 2) {
                    throw new Exception("Le 'target' doit etre au format 'Controller#Method'");
                }
                [$controller, $functionController] = $routeArray;
                if (isset($matches[1])) {
                    $instance = new $controller($matches[1]);
                    call_user_func_array([$instance, $functionController], $matches);
                    return;
                }
                (new $controller())->$functionController();
                return;
            }
        }
        http_response_code(404);
        page404View::show();
    }
}
\end{lstlisting}

\section{Jeux de tests complets}
\label{ann:tests}

Les tests sont écrits avec PHPUnit~11 et organisés en trois niveaux (unitaire, intégration,
sécurité). Ils se trouvent sous \file{tests/} et s'exécutent en une commande
(\code{php vendor/bin/phpunit}).

\begin{lstlisting}[style=text, caption={Arborescence des tests}]
tests/
  bootstrap.php             autoloader maison + Database
  UtilsTest.php             echappement XSS, AJAX, reponse JSON      (11)
  CsrfHelperTest.php        jeton CSRF : generation, verification     (8)
  SessionManagerTest.php    login / isAdmin au mock (DI, zero BDD)    (3)
  Integration/
    DatabaseTestCase.php    transaction + rollback par test
    UserModelTest.php       CRUD, hachage, UserValidator            (10)
    PostModelTest.php       creation, lecture, suppression           (6)
    SqlInjectionTest.php    charge d'injection inerte                (2)
\end{lstlisting}

\paragraph{Exemple — test d'intégration sécurité (injection SQL).}

\begin{lstlisting}[style=php, caption={\file{tests/Integration/SqlInjectionTest.php} (extrait)}]
final class SqlInjectionTest extends DatabaseTestCase
{
    // setUp() insere 1 utilisateur + 1 post au contenu unique (MARQUEUR),
    // le tout dans la transaction annulee au teardown.

    public function une_charge_d_injection_reste_inerte(): void
    {
        // "%' OR '1'='1%" est cherche litteralement : 0 resultat, sans exception.
        $r = (new SearchModel())->searchPosts("' OR '1'='1");
        $this->assertCount(0, $r);
    }

    public function une_recherche_normale_fonctionne(): void
    {
        $r = (new SearchModel())->searchPosts(self::MARQUEUR);
        $this->assertCount(1, $r);
    }
}
\end{lstlisting}

\paragraph{Sortie complète d'exécution.}

\begin{lstlisting}[style=text, caption={\code{php vendor/bin/phpunit --testdox}}]
Csrf Helper
 OK GetToken genere 64 caracteres hexadecimaux
 OK GetToken est stable dans une meme session
 OK VerifyToken accepte le bon token
 OK VerifyToken rejette un mauvais token
 OK VerifyToken rejette une chaine vide
 OK VerifyToken rejette null
 OK VerifyToken rejette quand aucun token en session
 OK Field rend un input cache avec le token
Post Model
 OK CreatePost insere et renvoie un id
 OK GetPostById lit le post cree
 OK GetPostById retourne null si inconnu
 OK Delete supprime le post
 OK Delete retourne false si post inexistant
 OK GetAllPostsByUserId ne retourne que les posts de l auteur
Session Manager
 OK Login peuple la session depuis le modele
 OK IsAdmin vrai pour le role admin
 OK IsAdmin faux pour un role non admin
Sql Injection
 OK Une charge d injection reste inerte
 OK Une recherche normale fonctionne
User Model
 OK GetUserByEmail retourne l utilisateur existant
 OK GetUserByEmail retourne null si inconnu
 OK Le mot de passe est stocke hashe
 OK CreateUser stocke un mot de passe hashe
 OK Login valide avec les bons identifiants
 OK Login refuse un mauvais mot de passe
 OK Login refuse un email inexistant
 OK Login refuse des champs vides
 OK Register refuse un email deja utilise
 OK Register accepte un email libre
Utils
 OK E neutralise une balise script
 OK E encode les chevrons
 OK E encode les guillemets doubles
 OK E encode l apostrophe en apos
 OK E encode l esperluette
 OK E transforme null en chaine vide
 OK IsAjax vrai avec l entete xmlhttprequest
 OK IsAjax faux sans l entete
 OK SendResponse serialise success et message
 OK SendResponse inclut data quand non vide
 OK SendResponse omet data quand falsy

OK (40 tests, 62 assertions)
\end{lstlisting}

\paragraph{Démonstrations de sécurité fonctionnelles.} Rejouées en \code{curl}, application
lancée ; preuve = statut HTTP + ligne de journal (le \emph{logger} écrit les avertissements sur
\code{stderr}, collectés par Docker).

\begin{lstlisting}[style=text, caption={CSRF (POST sans jeton) et IDOR (Bob supprime un message d'Alice)}]
$ curl -i -X POST http://localhost/login
HTTP/1.1 403 Forbidden
403 - Token CSRF refuse.
  journal: {"level":"warning","msg":"csrf.token.rejected","context":{"uri":"/login"}}

$ curl -i -X POST .../api/posts/delete   (session Bob, postId d'Alice = 1)
HTTP/1.1 403 Forbidden
{"success":false,"message":"Vous n'avez pas la permission de supprimer ce post."}
  journal: {"level":"warning","msg":"post.delete.forbidden",
           "context":{"current_user_id":2,"post_owner_id":1,"post_id":1}}
\end{lstlisting}

\paragraph{Jeu d'essai « Publier un message » — données obtenues (16 juin 2026).}

\begin{lstlisting}[style=text, caption={Réponses obtenues pour les sept cas}]
Nominal            -> success=true, id=53, HTTP 200
Contenu vide       -> success=false, "Donnees invalides ou contenu vide"
Image trop lourde  -> success=false, "Image trop lourde (max 2M)."
Extension .php     -> success=false, "Type de fichier non autorise."
Charge XSS         -> stocke brut "<script>...", rendu "&lt;script&gt;" (echappe)
Non authentifie    -> HTTP 302 -> /login   (AuthMiddleware)
Sans jeton CSRF    -> HTTP 403             (CsrfMiddleware)
\end{lstlisting}

\section{Maquettes et captures d'écran}
\label{ann:maquettes}

Maquettes des écrans de \projet{} réalisées sous Penpot (versions bureau et mobile), complétant
les trois écrans structurants présentés dans le corps du dossier (section~\ref{sec:maquettes}).

\subsection{Écrans bureau}

\begin{figure}[H]
\centering
\includegraphics[width=0.8\textwidth]{desktop-inscription.png}
\caption{Inscription (bureau) : même gabarit que la connexion, étendu aux champs du compte.}
\end{figure}

\begin{figure}[H]
\centering
\includegraphics[width=0.92\textwidth]{desktop-admin.png}
\caption{Vue bureau avec barre latérale de navigation persistante (accueil, recherche, profil,
administration) : recherche d'utilisateurs et de messages, et résultats assortis de l'action de
modération « Supprimer ».}
\end{figure}

\subsection{Écrans mobiles}

\begin{figure}[H]
\centering
\includegraphics[height=7.2cm]{mobile-connexion.png}\hfill
\includegraphics[height=7.2cm]{mobile-inscription.png}\hfill

\includegraphics[height=7.2cm]{mobile-profil.png}
\caption{Mobile : connexion, inscription et profil utilisateur.}
\end{figure}

\begin{figure}[H]
\centering

\includegraphics[height=7.2cm]{mobile-recherche.png}\hfill
\includegraphics[height=7.2cm]{mobile-admin.png}
\caption{Mobile : recherche et administration. La barre de navigation inférieure reste constante.}
\end{figure}

\paragraph{Code source complet.} Le code source complet de \projet{} (contrôleurs, modèles, vues,
middlewares, configuration Docker) est disponible dans le dépôt Git du projet. Les extraits
présentés dans ce dossier sont les plus significatifs au regard des compétences évaluées.
